\section{案例化路线图：12 周开放 Agent 研究冲刺}

\subsection{从课程理念到执行计划}
为了把本讲的理念转成可操作项目，可以采用“12 周冲刺”结构：前 4 周打基础设施，中间 4 周做系统迭代，最后 4 周做稳健性和复现封装。这个节奏对应讲者的核心精神：\textbf{先把研究系统做成可验证机器，再追求局部性能峰值。}

\begin{figure}[H]
\centering
\includegraphics[width=0.93\textwidth]{frame-02-benefits.jpg}
\caption{Agent 系统收益目标：接口、能力、架构三位一体}
\end{figure}
\noindent\footnotesize 画面证据：字幕约在 00:00:43--00:01:39，讲者从“access shapes research”延展到系统能力边界。\normalsize

\begin{importantbox}{建议的冲刺目标不是“上线一个酷炫 demo”}
更高价值的目标是：
\begin{itemize}
    \item 任一实验结论能被同伴在 24 小时内复现；
    \item 任一失败轨迹能被快速定位到具体环节；
    \item 任一提升都能说明来自数据、策略还是工具链。
\end{itemize}
这样做短期看慢，长期看会显著降低“伪进步”比例。
\end{importantbox}

\subsection{阶段里程碑与交付物}
\begin{table}[H]
\centering
\begin{tabular}{p{2.2cm}p{3.5cm}p{3.6cm}p{3.2cm}}
\toprule
\textbf{阶段} & \textbf{核心任务} & \textbf{强制交付物} & \textbf{验收标准} \\
\midrule
第 1--2 周 & 统一任务协议与日志格式 & task schema, run schema, replay script & 任意任务可回放 \\
第 3--4 周 & 基线 agent 与评测跑通 & baseline result, failure taxonomy & 有稳定 baseline 方差 \\
第 5--6 周 & 引入记忆/反思/规划模块 & ablation report, trace dashboard & 能解释增益来源 \\
第 7--8 周 & 多 agent 编排与权限治理 & orchestration spec, safety gate & 越权行为可拦截可审计 \\
第 9--10 周 & 鲁棒性与分布外压力测试 & stress suite, adversarial suite & 不同分布下失败可分解 \\
第 11--12 周 & 复现封装与文档化发布 & reproducibility package, model card & 第三方可独立复现 \\
\bottomrule
\end{tabular}
\caption{12 周冲刺的最小研究工程化路径}
\end{table}

\begin{knowledgebox}{为什么先做 replay 再做优化}
很多团队把 80\% 时间花在“想新策略”，却没有可重放轨迹。  
没有 replay，你无法判断改进是策略有效还是随机波动；没有稳定 replay，后续所有优化都无法累积为可信知识。
\end{knowledgebox}

\subsection{建议的实验主循环（伪代码）}
\begin{lstlisting}[caption={Research Loop: 可复现 Agent 迭代主循环},label={lst:research-loop}]
for task_batch in benchmark:
    run_id = launch_agent(task_batch, config, tools, safety_policy)
    traces = collect_traces(run_id)
    metrics = evaluate(traces, eval_protocol)

    if metrics.regression_detected:
        rollback(config)
        tag_failure(traces, taxonomy)
        continue

    insights = reflection_and_ablation(traces, metrics)
    config = update_config(config, insights)
    snapshot(run_id, config, metrics, traces)
\end{lstlisting}

\begin{warningbox}{最容易忽视的工程债务：不可比较的实验}
如果每次迭代都更换任务集合、工具版本或评测脚本，就会出现“每次都在进步，但没有两次结果可比较”的假象。  
冲刺阶段必须冻结关键协议，只在受控窗口调整变量。
\end{warningbox}

\subsection{本章小结}
这个 12 周路线图的目标是把“开放科学 + Agent 工程”的理念落地为可执行制度：可复现、可审计、可解释。只要三者缺一，团队就会在规模化阶段反复返工。


